\section{Sharp operator boundaries and the actual Walsh spectrum}
\label{sec:tpc45-boundaries}

The fiber loss in the preceding sections is an exact operator norm, not
an artifact of Cauchy--Schwarz.  We first record the finite-dimensional
statement.

Let \(E\) be a finite atom set, let \(\pi:E\to W\) be a label map, and
define
\begin{equation}
 (G_\pi x)(w):=\sum_{e:\pi(e)=w}x(e).
 \label{eq:tpc45-grouping-operator}
\end{equation}

\begin{proposition}[Sharp grouping norm]
\label{prop:tpc45-sharp-grouping-norm}
If
\begin{equation}
 \rho(\pi):=\max_w|\pi^{-1}(w)|,
 \label{eq:tpc45-grouping-multiplicity}
\end{equation}
then
\begin{equation}
 \boxed{
 \|G_\pi\|_{\ell^2(E)\to\ell^2(W)}^2=\rho(\pi).}
 \label{eq:tpc45-exact-grouping-norm}
\end{equation}
\end{proposition}

\begin{proof}
Cauchy--Schwarz on each fiber gives the upper bound.  On a fiber of
maximal cardinality, take \(x\) constant and nonzero and set it to zero
off that fiber.  The quotient of the squared norms is exactly the fiber
cardinality.
\end{proof}

For large-prime cofactor descent, all atoms in one descended fiber have
the same physical M\"obius phase.  Indeed, if \(s=dr\), then
\begin{equation}
 \mu(d)\mu(qr)=-\mu(dr)=-\mu(s).
 \label{eq:tpc45-fiber-physical-coherence}
\end{equation}
Thus M\"obius does not cancel different points of a fixed-determinant
ladder after their common large prime has been removed.  Any power
improvement over the support exponent in
\eqref{eq:tpc45-unbalanced-fiber-bound}, beyond possible logarithmic
thinning from primality, must use the actual complex source/mask
coefficients or cancellation between different residual labels \(s\);
it cannot arise from independent M\"obius signs inside one fiber.

\subsection{What is proved directly about
\texorpdfstring{\(C_X(\kappa)\)}{C-X(kappa)}}

Let \(C_X(\kappa)\) be the actual Walsh coefficient defined by
\begin{equation}
 Z_0(\varepsilon)
 =\cD_0+\sum_{\kappa>1}C_X(\kappa)\chi_\kappa(\varepsilon).
 \label{eq:tpc45-CX-definition}
\end{equation}
Every equality label contains at most one prime from
\(\cP_{\rm bal}\).  Therefore
\begin{equation}
 C_X(\kappa)=0
 \quad\text{if \(\kappa\) contains at least three primes from
 \(\cP_{\rm bal}\)}.
 \label{eq:tpc45-no-three-band-primes}
\end{equation}
There is also no pure singleton coefficient:
\begin{equation}
 C_X(q)=0
 \qquad(q\in\cP_{\rm bal}).
 \label{eq:tpc45-no-pure-singleton}
\end{equation}
Indeed, a pair kernel equal to \(q\) would require an untouched label
\(u\ge d_-N_-\) to equal a touched label with \(q\) deleted, which is
at most \(d_+\sqrt{N_+}\).  This contradicts
\eqref{eq:tpc45-touched-untouched-separation}.

For distinct \(p,q\in\cP_{\rm bal}\), a pure coefficient \(C_X(pq)\)
can occur only when the two labels delete to the same residual label
\(s\).  Introduce the Hilbert space with coordinates \((i,s)\), and
let \(x_q(i,s)\) be the coefficient of the unique row carrying \(q\)
and deleting to \(s\), with value zero when that row is absent.
Uniqueness follows because two such rows would have the same original
label \(qs\) in one coordinate, contrary to
\cref{lem:tpc45-label-injectivity}.  Then
\begin{align}
 C_X(pq)&=2\operatorname{Re}\langle x_p,x_q\rangle,
 \label{eq:tpc45-Cpq-Gram}\\
 \cD_{\rm bal}
 :=\sum_{q\in\cP_{\rm bal}}\|x_q\|^2
 &\le\cD_0.
 \label{eq:tpc45-band-atom-energy}
\end{align}
The second line is the atomic diagonal of the band-touched atoms:
equality-label injectivity makes each fixed \((i,s,q)\) entry a
singleton, and the same atom cannot carry two band primes.

Averaging every prime sign outside the band leaves the exact quadratic
projection
\begin{equation}
 \boxed{
 \E_{\varepsilon_a:a\notin\cP_{\rm bal}}Z_0(\varepsilon)
 =\cD_0+
  \sum_{\substack{p<q\\p,q\in\cP_{\rm bal}}}
  C_X(pq)\varepsilon_p\varepsilon_q.}
 \label{eq:tpc45-exact-band-quadratic}
\end{equation}
At the physical band point this is equivalently
\begin{equation}
 \E_{a\notin\cP_{\rm bal}}
 Z_0(-\one_{\cP_{\rm bal}},\varepsilon)
 =\cD_0-\cD_{\rm bal}
  +\left\|\sum_{q\in\cP_{\rm bal}}x_q\right\|^2.
 \label{eq:tpc45-exact-band-Gram-collapse}
\end{equation}
Combining this identity with
\eqref{eq:tpc45-band-frozen-mean} proves the actual-coefficient estimate
\begin{equation}
 \cD_0+
 \sum_{\substack{p<q\\p,q\in\cP_{\rm bal}}}C_X(pq)
 \le X^{o(1)}\cD_0.
 \label{eq:tpc45-pure-band-spectrum-bound}
\end{equation}
This controls a genuine polynomial band of the physical rough spectrum;
it is not an estimate for every kernel containing a band prime.

\subsection{The sharp continuation boundary}

If the upper endpoint of \(\cP_{\rm bal}\) is crossed, cofactors
\(r<L_0\) appear and the exact norm in
\cref{prop:tpc45-sharp-grouping-norm} may cost \(L_0/r\).  If a physical
fiber contains \(H\) atoms, coefficients of equal phase supported on
that fiber give grouped energy \(H\) times their atomic energy.  This is
a sharp statement conditional on the declared support; it neither
asserts that the actual masks choose equal phases nor proves that a
fiber of polynomial cardinality occurs.

Accordingly, the results above do not show that the residual TPC--44
correlation product is small for all high primes, and they do not show
that it is large.  They prove a physical freezing theorem on the
balanced band and identify the first adjacent range where support and
diagonal information cease to suffice.  No fixed-shift Chowla estimate,
prime-pair asymptotic, parity breach, or twin-prime conclusion follows.
